\section{Design principles, limitations, and outlook}
\label{sec:outlook}

The engine is deliberately small, and we intend to keep it so while widening what
it can certify. This closing section distils the principles that the preceding
sections established, states the limitations plainly, and describes the
self-renewing programme we are building towards. The throughline is unchanged
from the first version of the engine: a research result should be a re-runnable,
self-checking object, and the infrastructure that makes it so is worth building
once, carefully, and maintaining in the open.

\subsection{The design principles, distilled}

Five principles carry the design. Each was learned in operation, several the hard
way, and each is stated here in the form the project applies to itself.

\begin{principle}[Parameterised-only registry]
A request names a method and schema-validated parameters; it never supplies code.
The registry is therefore the primary guard against arbitrary-code execution, and
the operating-system sandbox is defence in depth on top of it, not the
load-bearing control.
\end{principle}

\begin{principle}[A failable test per method]
A method is not in the catalogue until it ships with a pre-registered, failable
test of its own output. Holes in coverage are marked, never filled. The registered
preference is explicit: we would rather have twenty-six checked methods than
fifty asserted ones.
\end{principle}

\begin{principle}[Honest negatives are results]
A result that does not clear its test, or that contradicts an earlier finding, is
reported as such and treated as an immutable fact of the record. The
not-significant SOCH-C test ($p=0.105$), the degenerate false-discovery-controlled
NAMH network ($0/552$), and the below-chance MCPFM v2 re-evaluation (AUC $0.470$)
are findings, not blemishes to be managed away.
\end{principle}

The remaining two principles are properties the system maintains rather than
preferences it expresses, so we state them as invariants.

\begin{invariant}[Governed determinism]
Worker count never changes a number. A single governed core budget, clamped last
by a server-set hard ceiling, with linear-algebra threading pinned to one thread
per spawned process, makes every result independent of how many workers run it.
The transfer-entropy loops are order-free, so this is determinism-safe, and the
GPU offload of the same computation is held bit-exact to the CPU estimator on the
only gated quantity, the observed transfer entropy, to within machine epsilon.
\end{invariant}

\begin{invariant}[Machine-written ledgers are never hand-edited]
The single evaluation artefact, the state ledger's measured block, and the claims
record are written by their generators and never by hand. A regression turns the
public dashboard red rather than passing silently; the state ledger is loud about
drift; and the claims record never deletes, so a result that turns red becomes a
contested pair rather than a quiet disappearance.
\end{invariant}

These five are not independent. The failable-test principle is what makes the
honest-negative principle enforceable; the never-hand-edit invariant is what
makes both credible to a reader who was not in the room. Together they are the
sense in which the engine's reliability is a property of its process rather than a
promise from its authors.

\subsection{Limitations}

We state the limitations plainly, because the discipline that reports honest
negatives in the elements applies equally to the engine that carries them.

\paragraph{Reported degeneracies and negatives.} Some of the programme's
headline tests do not clear, and we have not hidden this. The
false-discovery-controlled NAMH network is degenerate, retaining $0$ of $552$
candidate edges; the engine's NAMH centralities are magnitude-ordered and are not
a claim that the FDR-gated network is non-empty. The pre-registered MCPFM v2
re-evaluation returns an area under the ROC curve of $0.470$, below chance. The
SOCH-C significance test is not significant at $p=0.105$. None of these is a bug
to be fixed; each is a fact of the record.

\paragraph{Infrastructure constraints.} The heavy asynchronous estimators run on
a single workstation tower, so the engine's throughput on the four tower methods
is bounded by one machine and its one GPU; a second compute node is a
pre-registered future step, not a present capability. The natural-language layer
is rate-limited and budget-capped: over its daily cap it returns an explicit
capacity error with a reset hint, never a fabricated answer, which is the right
failure but is a hard ceiling on interactive use. The formal layer is partial:
of twelve Lean~4 declarations, eleven are machine-accepted and one, the analytic
half of the peak-location lemma, honestly carries a \code{sorry}; the spectral
theory is closed where the proposition $\Norm{H(\omega)}^2 = S(\omega)$ is
concerned, but the peak-location chain is not yet complete.

\paragraph{Scope.} The engine certifies what its twenty-six methods cover and no
more. A question outside the catalogue has no banded answer, by design, and a
figure outside the verified record is marked for verification rather than
asserted. This is a deliberate narrowness, but it is a narrowness, and the
honest description of the engine is that it makes a fixed catalogue
re-runnable, not that it makes financial econometrics in general reproducible.

\subsection{Towards a self-renewing programme}

The direction of travel is to make the engine adaptive in the same disciplined
sense the markets it studies are adaptive: a system that learns from its own
operation without ever loosening the gate that keeps it honest.

\paragraph{The nightly loop as a learning system.} The engine's daily systemic-risk
tick already runs unattended on a managed scheduler. A reboot-surviving nightly loop,
now installed in cron, re-arms the tower, runs the
full public suite as one genuine execution with the asynchronous rows submitted
as real jobs, refreshes the state ledger and appends any drift, refreshes the
claims record without ever deleting, and renders and publishes an academic report
of the night's run. Crucially, the loop's own wiring is itself a failable
evaluation: the system that checks the engine is checked by the same kind of
pre-registered test it applies to everything else. The next step is to treat the
sequence of nightly runs as a learning record, so that drift over time becomes a
first-class signal rather than a per-night log line.

\paragraph{Pre-registered, PI-gated gap proposals.} A gap-proposal pipeline reads
the engine's own usage and coverage and proposes candidate additions to the
catalogue. It is pre-registered and never auto-deploys: a proposal is a review
document for the principal investigator, not a change to the live registry. When
the pipeline ran end-to-end and its rule found no qualifying gap, that
no-proposal outcome was recorded as a first-class result, in the same spirit as a
not-significant test. The design intent is a catalogue that can grow on evidence
while every addition still passes through the same human gate and the same
failable-test requirement.

\paragraph{More elements, more formal coverage.} Two expansions are in view on the
same terms. First, more elements may join the programme, each only with a failable
evaluation and each instantiating the shared substrate, so that the marginal cost
of a new framework is the cost of its configuration and its test, not a new
codebase. Second, the formal layer can grow: closing the remaining \code{sorry} in
the SOCH peak-location chain, and extending machine-checked proof to results
beyond SOCH, would widen the set of claims the engine can certify at the level of
a checked proof rather than a banded number.

\paragraph{The self-renewing vision.} Taken together, these point at a programme
that renews itself: a fixed, sandboxed engine; a shared substrate that new
elements configure rather than reinvent; a nightly loop that learns from its own
record; and a gap pipeline that proposes growth under a standing human gate. We
do not claim this is finished, and we have stated where it is partial. We claim
only that the parts are built, that they hold to the standard we have described,
and that the standard, not any single number, is the contribution we are most
willing to defend.
